\documentclass[10pt]{article}
\usepackage[utf8]{inputenc}
\usepackage[T1]{fontenc}
\usepackage{lmodern}
\usepackage[paperwidth=384pt,paperheight=900pt,margin=10pt]{geometry}
\usepackage{amsmath,amssymb}
\usepackage{array}
\usepackage{enumitem}
\usepackage{xcolor}
\usepackage{microtype}

\setlength{\parindent}{0pt}
\setlength{\parskip}{4pt}
\setlist[itemize]{leftmargin=12pt}
\setlist[enumerate]{leftmargin=14pt}

\sloppy
\allowdisplaybreaks
\emergencystretch=2em

\begin{document}

\section*{Uppgift 1}

\subsection*{Original question}

För händelserna \(A\), \(B\) och \(C\) är följande
sannolikheter kända:

\[
P(A)=\frac{1}{3},
\qquad
P(B)=\frac{1}{4},
\qquad
P(C)=\frac{1}{2}.
\]

Dessutom vet man att
\[
P(A\cup B\cup C)=1
\]
och att
\[
P(A\cap C)=P(B\cap C)=0.
\]

\begin{enumerate}[label=\alph*)]
\item Bevisa att
\[
P(A\cup B)=\frac{1}{2}.
\]
\hfill \((2.5p)\)

\item Bestäm \(P(A\cap B)\) och \(P(A\mid B)\).
\hfill \((2p)\)

\item Vilka par av händelserna \(A\), \(B\), \(C\)
är oberoende? Motivera dina svar!
\hfill \((1.5p)\)
\end{enumerate}

\subsection*{Compiled notes from Yacoub + Obid}
% UNCERTAIN: The supplied note images are not visibly labelled by author,
% so Yacoub's and Obid's notes are compiled together as the provided notes.

\subsubsection*{What we are doing}

We use probability rules for unions,
intersections, conditional probability, and
independence.

The notes notice that:
\[
P(A\cap C)=0,
\qquad
P(B\cap C)=0.
\]

So \(A\) and \(C\) are mutually exclusive, and
\(B\) and \(C\) are mutually exclusive.

\subsubsection*{Given information}

\[
P(A)=\frac{1}{3}.
\]

\[
P(B)=\frac{1}{4}.
\]

\[
P(C)=\frac{1}{2}.
\]

\[
P(A\cup B\cup C)=1.
\]

\[
P(A\cap C)=0.
\]

\[
P(B\cap C)=0.
\]

\subsubsection*{Step-by-step solution}

\textbf{a) Show that \(P(A\cup B)=\frac12\)}

The notes use that \(C\) does not overlap with
\(A\) or \(B\).

So \(C\) does not overlap with \(A\cup B\):
\[
(A\cup B)\cap C=\varnothing
\]
in probability.

Therefore:
\[
P(A\cup B\cup C)
=
P(A\cup B)+P(C).
\]

Now insert the known values:
\begin{align*}
1
&=
P(A\cup B)
+
\frac{1}{2}
\\[2mm]
P(A\cup B)
&=
1-\frac{1}{2}
\\[2mm]
&=
\frac{1}{2}.
\end{align*}

So:
\[
P(A\cup B)=\frac{1}{2}.
\]

\textbf{Extra formula shown in the notes}

The notes also write the full three-event union
formula:
\begin{align*}
P(A\cup B\cup C)
&=
P(A)+P(B)+P(C)
\\
&\quad
-P(A\cap B)
-P(A\cap C)
\\
&\quad
-P(B\cap C)
+P(A\cap B\cap C).
\end{align*}

For two events, the notes use:
\[
P(A\cup B)
=
P(A)+P(B)-P(A\cap B).
\]

\textbf{b) Find \(P(A\cap B)\)}

För händelserna \(A\), \(B\) och \(C\) är följande
sannolikheter kända:

\[
P(A)=\frac{1}{3},
\qquad
P(B)=\frac{1}{4},
\qquad
P(C)=\frac{1}{2}.
\]

Dessutom vet man att
\[
P(A\cup B\cup C)=1
\]
och att
\[
P(A\cap C)=P(B\cap C)=0.
\]

\begin{enumerate}[label=\alph*)]
\item Bevisa att
\[
P(A\cup B)=\frac{1}{2}.
\]
\hfill \((2.5p)\)

\item Bestäm \(P(A\cap B)\) och \(P(A\mid B)\).
\hfill \((2p)\)

\item Vilka par av händelserna \(A\), \(B\), \(C\)
är oberoende? Motivera dina svar!
\hfill \((1.5p)\)
\end{enumerate}

Use the two-event union formula:
\[
P(A\cup B)
=
P(A)+P(B)-P(A\cap B).
\]

Insert values from part a) and the question:
\begin{align*}
\frac{1}{2}
&=
\frac{1}{3}
+
\frac{1}{4}
-
P(A\cap B).
\end{align*}

Move terms around:
\begin{align*}
P(A\cap B)
&=
\frac{1}{3}
+
\frac{1}{4}
-
\frac{1}{2}
\\[2mm]
&=
\frac{4}{12}
+
\frac{3}{12}
-
\frac{6}{12}
\\[2mm]
&=
\frac{1}{12}.
\end{align*}

So:
\[
P(A\cap B)=\frac{1}{12}.
\]

\textbf{b) Find \(P(A\mid B)\)}

För händelserna \(A\), \(B\) och \(C\) är följande
sannolikheter kända:

\[
P(A)=\frac{1}{3},
\qquad
P(B)=\frac{1}{4},
\qquad
P(C)=\frac{1}{2}.
\]

Dessutom vet man att
\[
P(A\cup B\cup C)=1
\]
och att
\[
P(A\cap C)=P(B\cap C)=0.
\]

\begin{enumerate}[label=\alph*)]
\item Bevisa att
\[
P(A\cup B)=\frac{1}{2}.
\]
\hfill \((2.5p)\)

\item Bestäm \(P(A\cap B)\) och \(P(A\mid B)\).
\hfill \((2p)\)

\item Vilka par av händelserna \(A\), \(B\), \(C\)
är oberoende? Motivera dina svar!
\hfill \((1.5p)\)
\end{enumerate}

The notes use the conditional probability formula:
\[
P(A\mid B)
=
\frac{P(A\cap B)}{P(B)}.
\]

Insert the values:
\begin{align*}
P(A\mid B)
&=
\frac{\frac{1}{12}}{\frac{1}{4}}
\\[2mm]
&=
\frac{1}{12}\cdot 4
\\[2mm]
&=
\frac{1}{3}.
\end{align*}

So:
\[
P(A\mid B)=\frac{1}{3}.
\]

\textbf{c) Which pairs are independent?}

För händelserna \(A\), \(B\) och \(C\) är följande
sannolikheter kända:

\[
P(A)=\frac{1}{3},
\qquad
P(B)=\frac{1}{4},
\qquad
P(C)=\frac{1}{2}.
\]

Dessutom vet man att
\[
P(A\cup B\cup C)=1
\]
och att
\[
P(A\cap C)=P(B\cap C)=0.
\]

\begin{enumerate}[label=\alph*)]
\item Bevisa att
\[
P(A\cup B)=\frac{1}{2}.
\]
\hfill \((2.5p)\)

\item Bestäm \(P(A\cap B)\) och \(P(A\mid B)\).
\hfill \((2p)\)

\item Vilka par av händelserna \(A\), \(B\), \(C\)
är oberoende? Motivera dina svar!
\hfill \((1.5p)\)
\end{enumerate}

The notes use this test:
\[
P(X\cap Y)
=
P(X)\cdot P(Y).
\]

If both sides are equal, the two events are
independent.

\textbf{Test 1: \(A\) and \(B\)}

From part b):
\[
P(A\cap B)=\frac{1}{12}.
\]

Now calculate:
\begin{align*}
P(A)\cdot P(B)
&=
\frac{1}{3}\cdot\frac{1}{4}
\\[2mm]
&=
\frac{1}{12}.
\end{align*}

They are equal:
\[
\frac{1}{12}=\frac{1}{12}.
\]

So:
\[
A \text{ and } B
\text{ are independent.}
\]

\textbf{Test 2: \(A\) and \(C\)}

From the question:
\[
P(A\cap C)=0.
\]

Now calculate:
\begin{align*}
P(A)\cdot P(C)
&=
\frac{1}{3}\cdot\frac{1}{2}
\\[2mm]
&=
\frac{1}{6}.
\end{align*}

They are not equal:
\[
0\ne \frac{1}{6}.
\]

So:
\[
A \text{ and } C
\text{ are not independent.}
\]

\textbf{English clarification:}
The handwritten note labels this product like
\(P(A)\cdot P(B)\), but the numbers used are
\(\frac13\cdot\frac12\). That is the check for
\(A\) and \(C\).

\textbf{Test 3: \(B\) and \(C\)}

From the question:
\[
P(B\cap C)=0.
\]

Now calculate:
\begin{align*}
P(B)\cdot P(C)
&=
\frac{1}{4}\cdot\frac{1}{2}
\\[2mm]
&=
\frac{1}{8}.
\end{align*}

They are not equal:
\[
0\ne \frac{1}{8}.
\]

So:
\[
B \text{ and } C
\text{ are not independent.}
\]

\subsubsection*{Final answer}

\begin{enumerate}[label=\alph*)]
\item
\[
P(A\cup B)=\frac{1}{2}.
\]

\item
\[
P(A\cap B)=\frac{1}{12}.
\]

\[
P(A\mid B)=\frac{1}{3}.
\]

\item
\[
A \text{ and } B
\text{ are independent.}
\]

\[
A \text{ and } C
\text{ are not independent.}
\]

\[
B \text{ and } C
\text{ are not independent.}
\]
\end{enumerate}

\subsubsection*{What to remember}

For two events:
\[
P(A\cup B)
=
P(A)+P(B)-P(A\cap B).
\]

Conditional probability:
\[
P(A\mid B)
=
\frac{P(A\cap B)}{P(B)}.
\]

Independence test:
\[
P(X\cap Y)
=
P(X)P(Y).
\]

\section*{Uppgift 2}

\subsection*{Original question}

Ett spel börjar med ett tärningskast med resultatet
\[
k\in \{1,2,3,4,5,6\}.
\]

Om \(k\) är udda vinner spelaren \(k\) kronor.
Om \(k\) är jämnt kastar han en gång till.
Om andra kastet \(\ell\) är lika med 1 vinner han
\(k\) kronor. Annars förlorar han \(k\) kronor.

\begin{enumerate}[label=\alph*)]
\item Visualisera alla spelförlopp genom ett
träddiagram. \((2p)\)

\item Betrakta den stokastiska variabeln \(X\) som
beskriver spelets vinst (med negativa värden för
förluster). Bestäm sannolikhetsfunktionen för \(X\).
\((2p)\)

\item Beräkna \(E(X)\) och tolka resultatet ur
perspektivet för en spelare som upprepar spelet
många gånger. \((2p)\)
\end{enumerate}

\subsection*{Compiled notes from Yacoub + Obid}
% UNCERTAIN: The supplied note images are not visibly labelled by author,
% so Yacoub's and Obid's notes are compiled together as the provided notes.

\subsubsection*{What we are doing}

The game starts with one die roll.
Call the first result \(k\).

The notes split the game into:
\begin{itemize}
\item odd \(k\): win directly,
\item even \(k\): roll one more time,
\item second roll \(\ell=1\): win,
\item second roll \(\ell\ne 1\): lose.
\end{itemize}

The random variable \(X\) is the final profit.
Wins are positive and losses are negative.

\subsubsection*{Given information}

\[
k\in\{1,2,3,4,5,6\}.
\]

Odd first rolls:
\[
k\in\{1,3,5\}.
\]

Even first rolls:
\[
k\in\{2,4,6\}.
\]

For each die roll:
\[
P(\text{specific number})=\frac{1}{6}.
\]

For the second roll:
\[
P(\ell=1)=\frac{1}{6},
\qquad
P(\ell\ne 1)=\frac{5}{6}.
\]

\subsubsection*{Step-by-step solution}

\textbf{a) Tree diagram}

The notes draw the game tree. Written in narrow form,
the same tree is:

\begin{itemize}
\item
First roll \(k=1\), probability \(\frac16\):
\[
X=1
\quad\text{win}.
\]

\item
First roll \(k=3\), probability \(\frac16\):
\[
X=3
\quad\text{win}.
\]

\item
First roll \(k=5\), probability \(\frac16\):
\[
X=5
\quad\text{win}.
\]

\item
First roll \(k=2\), probability \(\frac16\):
\[
\ell=1
\text{ with probability }\frac16
\Rightarrow X=2.
\]
\[
\ell\ne 1
\text{ with probability }\frac56
\Rightarrow X=-2.
\]

\item
First roll \(k=4\), probability \(\frac16\):
\[
\ell=1
\text{ with probability }\frac16
\Rightarrow X=4.
\]
\[
\ell\ne 1
\text{ with probability }\frac56
\Rightarrow X=-4.
\]

\item
First roll \(k=6\), probability \(\frac16\):
\[
\ell=1
\text{ with probability }\frac16
\Rightarrow X=6.
\]
\[
\ell\ne 1
\text{ with probability }\frac56
\Rightarrow X=-6.
\]
\end{itemize}

\textbf{b) Probability mass function for \(X\)}

Ett spel börjar med ett tärningskast med resultatet
\[
k\in \{1,2,3,4,5,6\}.
\]

Om \(k\) är udda vinner spelaren \(k\) kronor.
Om \(k\) är jämnt kastar han en gång till.
Om andra kastet \(\ell\) är lika med 1 vinner han
\(k\) kronor. Annars förlorar han \(k\) kronor.

\begin{enumerate}[label=\alph*)]
\item Visualisera alla spelförlopp genom ett
träddiagram. \((2p)\)

\item Betrakta den stokastiska variabeln \(X\) som
beskriver spelets vinst (med negativa värden för
förluster). Bestäm sannolikhetsfunktionen för \(X\).
\((2p)\)

\item Beräkna \(E(X)\) och tolka resultatet ur
perspektivet för en spelare som upprepar spelet
många gånger. \((2p)\)
\end{enumerate}

The notes say that \(X\) represents the final
winnings.

For odd first rolls \(1,3,5\), the player wins
directly. So:
\[
P(X=1)=\frac16=\frac{6}{36},
\]
\[
P(X=3)=\frac16=\frac{6}{36},
\]
\[
P(X=5)=\frac16=\frac{6}{36}.
\]

For positive even wins \(2,4,6\), we need:
\[
\text{first roll is }k
\quad\text{and}\quad
\ell=1.
\]

So each such probability is:
\[
\frac16\cdot\frac16=\frac{1}{36}.
\]

Therefore:
\[
P(X=2)=\frac{1}{36},
\]
\[
P(X=4)=\frac{1}{36},
\]
\[
P(X=6)=\frac{1}{36}.
\]

For losses \(-2,-4,-6\), we need:
\[
\text{first roll is }k
\quad\text{and}\quad
\ell\ne 1.
\]

So each such probability is:
\[
\frac16\cdot\frac56=\frac{5}{36}.
\]

Therefore:
\[
P(X=-2)=\frac{5}{36},
\]
\[
P(X=-4)=\frac{5}{36},
\]
\[
P(X=-6)=\frac{5}{36}.
\]

\textbf{Sanity check from the notes}

Add all probabilities:
\begin{align*}
&\frac{6+6+6+1+1+1+5+5+5}{36}
\\[2mm]
&=
\frac{36}{36}
\\[2mm]
&=1.
\end{align*}

So the probability mass function makes sense.

\textbf{c) Calculate \(E(X)\)}

Ett spel börjar med ett tärningskast med resultatet
\[
k\in \{1,2,3,4,5,6\}.
\]

Om \(k\) är udda vinner spelaren \(k\) kronor.
Om \(k\) är jämnt kastar han en gång till.
Om andra kastet \(\ell\) är lika med 1 vinner han
\(k\) kronor. Annars förlorar han \(k\) kronor.

\begin{enumerate}[label=\alph*)]
\item Visualisera alla spelförlopp genom ett
träddiagram. \((2p)\)

\item Betrakta den stokastiska variabeln \(X\) som
beskriver spelets vinst (med negativa värden för
förluster). Bestäm sannolikhetsfunktionen för \(X\).
\((2p)\)

\item Beräkna \(E(X)\) och tolka resultatet ur
perspektivet för en spelare som upprepar spelet
många gånger. \((2p)\)
\end{enumerate}

The notes use:
\[
E(X)=\sum xP(X=x).
\]

Put in all possible values:
\begin{align*}
E(X)
&=
1\left(\frac{6}{36}\right)
+3\left(\frac{6}{36}\right)
+5\left(\frac{6}{36}\right)
\\
&\quad
+2\left(\frac{1}{36}\right)
+4\left(\frac{1}{36}\right)
+6\left(\frac{1}{36}\right)
\\
&\quad
-2\left(\frac{5}{36}\right)
-4\left(\frac{5}{36}\right)
-6\left(\frac{5}{36}\right).
\end{align*}

Now collect the positive and negative parts:
\begin{align*}
E(X)
&=
\frac{66-60}{36}
\\[2mm]
&=
\frac{6}{36}
\\[2mm]
&=
\frac{1}{6}.
\end{align*}

So:
\[
E(X)=\frac{1}{6}.
\]

The notes interpret this as:
the player statistically makes
\[
\frac{1}{6}
\]
krona per game in the long run.

\subsubsection*{Final answer}

\begin{enumerate}[label=\alph*)]
\item
The tree has direct wins for:
\[
k=1,3,5.
\]
For:
\[
k=2,4,6,
\]
there is a second roll:
\[
\ell=1 \Rightarrow \text{win }k,
\]
\[
\ell\ne 1 \Rightarrow \text{lose }k.
\]

\item
\[
P(X=1)=P(X=3)=P(X=5)=\frac{6}{36}.
\]

\[
P(X=2)=P(X=4)=P(X=6)=\frac{1}{36}.
\]

\[
P(X=-2)=P(X=-4)=P(X=-6)=\frac{5}{36}.
\]

\item
\[
E(X)=\frac{1}{6}.
\]

In the long run, the player wins on average
\[
\frac{1}{6}
\]
krona per game.
\end{enumerate}

\subsubsection*{What to remember}

For a discrete random variable:
\[
E(X)=\sum xP(X=x).
\]

For a path in a tree diagram, multiply the
probabilities along the path.

Always check that the probabilities add to:
\[
1.
\]

\section*{Uppgift 3}

\subsection*{Original question}

Betrakta funktionen

\[
f(x)=
\begin{cases}
-x+\sqrt{3},
& 0\le x\le a,\\
0,
& \text{annars}.
\end{cases}
\]

\begin{enumerate}[label=\alph*)]
\item Bestäm konstanten \(a\ge 0\) sådan att
\(f\) är en täthetsfunktion för en stokastisk
variabel \(X\). \((2p)\)

\item Beräkna
\[
p=P\left(X\le \frac{1}{\sqrt{3}}\right).
\]
\((1.5p)\)

\item Antag att svaret i b) är \(p=0,8\)
(vad man faktiskt får efter en grov avrundning).
Om vi genomför 8 oberoende försök
(som modelleras genom \(X\)) vad är sannolikheten
att vi får ett resultat mindre än eller lika med
\(\frac{1}{\sqrt{3}}\) fler än två gånger?
\((2.5p)\)

\item Fortsättning av c). Om vi fördubblar antalet
av våra försök vad är då sannolikheten att vi inte
får ett resultat mindre än eller lika med
\(\frac{1}{\sqrt{3}}\) i vart och ett av försöken?
\((2p)\)
\end{enumerate}

\subsection*{Compiled notes from Yacoub + Obid}
% UNCERTAIN: The supplied note images are not visibly labelled by author,
% so Yacoub's and Obid's notes are compiled together as the provided notes.

\subsubsection*{What we are doing}

First we choose \(a\) so that the total area under
the density is \(1\).

Then we calculate the probability in part b) by
integrating the density.

For parts c) and d), the notes use the binomial
distribution with success probability \(p=0.8\).

\subsubsection*{Given information}

\[
f(x)=
\begin{cases}
-x+\sqrt{3},
& 0\le x\le a,\\
0,
& \text{annars}.
\end{cases}
\]

\[
a\ge 0.
\]

For a density function:
\[
\int_{-\infty}^{\infty} f(x)\,dx=1.
\]

In parts c) and d), use:
\[
p=0.8.
\]

\subsubsection*{Step-by-step solution}

\textbf{a) Find \(a\)}

Since \(f(x)=0\) outside \(0\le x\le a\), the notes
write:
\[
\int_0^a
(-x+\sqrt{3})\,dx
=1.
\]

Integrate:
\begin{align*}
\int_0^a
(-x+\sqrt{3})\,dx
&=
\left[
-\frac{x^2}{2}
+x\sqrt{3}
\right]_0^a
\\[2mm]
&=
-\frac{a^2}{2}
+a\sqrt{3}.
\end{align*}

Set this equal to \(1\):
\begin{align*}
-\frac{a^2}{2}
+a\sqrt{3}
&=1.
\end{align*}

Multiply by \(2\) to remove the fraction:
\begin{align*}
-a^2+2a\sqrt{3}
&=2.
\end{align*}

Rearrange:
\begin{align*}
a^2-2a\sqrt{3}+2
&=0.
\end{align*}

Use the quadratic formula:
\begin{align*}
a
&=
\sqrt{3}
\pm
\sqrt{(\sqrt{3})^2-2}
\\[2mm]
&=
\sqrt{3}
\pm
\sqrt{3-2}
\\[2mm]
&=
\sqrt{3}\pm 1.
\end{align*}

So the two algebraic candidates are:
\[
a=\sqrt{3}+1
\quad\text{or}\quad
a=\sqrt{3}-1.
\]

\textbf{English clarification:}
The density must not become negative on its support.
Since
\[
-x+\sqrt{3}\ge 0
\]
requires
\[
x\le \sqrt{3},
\]
we choose the smaller value:
\[
a=\sqrt{3}-1.
\]

\textbf{b) Calculate}
\[
p=P\left(X\le \frac{1}{\sqrt{3}}\right).
\]

Betrakta funktionen

\[
f(x)=
\begin{cases}
-x+\sqrt{3},
& 0\le x\le a,\\
0,
& \text{annars}.
\end{cases}
\]

\begin{enumerate}[label=\alph*)]
\item Bestäm konstanten \(a\ge 0\) sådan att
\(f\) är en täthetsfunktion för en stokastisk
variabel \(X\). \((2p)\)

\item Beräkna
\[
p=P\left(X\le \frac{1}{\sqrt{3}}\right).
\]
\((1.5p)\)

\item Antag att svaret i b) är \(p=0,8\)
(vad man faktiskt får efter en grov avrundning).
Om vi genomför 8 oberoende försök
(som modelleras genom \(X\)) vad är sannolikheten
att vi får ett resultat mindre än eller lika med
\(\frac{1}{\sqrt{3}}\) fler än två gånger?
\((2.5p)\)

\item Fortsättning av c). Om vi fördubblar antalet
av våra försök vad är då sannolikheten att vi inte
får ett resultat mindre än eller lika med
\(\frac{1}{\sqrt{3}}\) i vart och ett av försöken?
\((2p)\)
\end{enumerate}

The notes integrate from \(0\) to
\(\frac{1}{\sqrt{3}}\):
\begin{align*}
P\left(X\le \frac{1}{\sqrt{3}}\right)
&=
\int_0^{1/\sqrt{3}}
(-x+\sqrt{3})\,dx.
\end{align*}

Use the same antiderivative as in part a):
\begin{align*}
P\left(X\le \frac{1}{\sqrt{3}}\right)
&=
\left[
-\frac{x^2}{2}
+x\sqrt{3}
\right]_0^{1/\sqrt{3}}.
\end{align*}

Substitute the upper limit:
\begin{align*}
p
&=
-
\frac{(1/\sqrt{3})^2}{2}
+
\frac{1}{\sqrt{3}}\sqrt{3}
\\[2mm]
&=
-
\frac{1/3}{2}
+1
\\[2mm]
&=
-\frac{1}{6}
+1
\\[2mm]
&=
\frac{5}{6}.
\end{align*}

So:
\[
p=\frac{5}{6}.
\]

\textbf{c) More than two successes in 8 trials}

Betrakta funktionen

\[
f(x)=
\begin{cases}
-x+\sqrt{3},
& 0\le x\le a,\\
0,
& \text{annars}.
\end{cases}
\]

\begin{enumerate}[label=\alph*)]
\item Bestäm konstanten \(a\ge 0\) sådan att
\(f\) är en täthetsfunktion för en stokastisk
variabel \(X\). \((2p)\)

\item Beräkna
\[
p=P\left(X\le \frac{1}{\sqrt{3}}\right).
\]
\((1.5p)\)

\item Antag att svaret i b) är \(p=0,8\)
(vad man faktiskt får efter en grov avrundning).
Om vi genomför 8 oberoende försök
(som modelleras genom \(X\)) vad är sannolikheten
att vi får ett resultat mindre än eller lika med
\(\frac{1}{\sqrt{3}}\) fler än två gånger?
\((2.5p)\)

\item Fortsättning av c). Om vi fördubblar antalet
av våra försök vad är då sannolikheten att vi inte
får ett resultat mindre än eller lika med
\(\frac{1}{\sqrt{3}}\) i vart och ett av försöken?
\((2p)\)
\end{enumerate}

The question says to use:
\[
p=0.8.
\]

Let \(Y\) be the number of successes, where success
means:
\[
X\le \frac{1}{\sqrt{3}}.
\]

Then:
\[
Y\sim \operatorname{Bin}(8,0.8).
\]

We want:
\[
P(Y>2).
\]

The notes use the complement:
\[
P(Y>2)=1-P(Y\le 2).
\]

So:
\[
P(Y>2)
=
1-
\bigl(
P(Y=0)+P(Y=1)+P(Y=2)
\bigr).
\]

Use the binomial formula:
\[
P(Y=k)
=
\binom{n}{k}
p^k(1-p)^{n-k}.
\]

Calculate the three terms:
\begin{align*}
P(Y=0)
&=
\binom{8}{0}
(0.8)^0(0.2)^8
\\[2mm]
&=
0.00000256.
\end{align*}

\begin{align*}
P(Y=1)
&=
\binom{8}{1}
(0.8)^1(0.2)^7
\\[2mm]
&=
0.00008192.
\end{align*}

\begin{align*}
P(Y=2)
&=
\binom{8}{2}
(0.8)^2(0.2)^6
\\[2mm]
&=
0.00114688.
\end{align*}

Add them:
\begin{align*}
P(Y\le 2)
&=
0.00000256
+0.00008192
\\
&\quad
+0.00114688
\\[2mm]
&=
0.00123136.
\end{align*}

Subtract from \(1\):
\begin{align*}
P(Y>2)
&=
1-0.00123136
\\[2mm]
&=
0.99876864.
\end{align*}

So the probability is approximately:
\[
0.9988.
\]

\textbf{d) Double the number of trials}

Betrakta funktionen

\[
f(x)=
\begin{cases}
-x+\sqrt{3},
& 0\le x\le a,\\
0,
& \text{annars}.
\end{cases}
\]

\begin{enumerate}[label=\alph*)]
\item Bestäm konstanten \(a\ge 0\) sådan att
\(f\) är en täthetsfunktion för en stokastisk
variabel \(X\). \((2p)\)

\item Beräkna
\[
p=P\left(X\le \frac{1}{\sqrt{3}}\right).
\]
\((1.5p)\)

\item Antag att svaret i b) är \(p=0,8\)
(vad man faktiskt får efter en grov avrundning).
Om vi genomför 8 oberoende försök
(som modelleras genom \(X\)) vad är sannolikheten
att vi får ett resultat mindre än eller lika med
\(\frac{1}{\sqrt{3}}\) fler än två gånger?
\((2.5p)\)

\item Fortsättning av c). Om vi fördubblar antalet
av våra försök vad är då sannolikheten att vi inte
får ett resultat mindre än eller lika med
\(\frac{1}{\sqrt{3}}\) i vart och ett av försöken?
\((2p)\)
\end{enumerate}

We double the number of trials:
\[
n=16.
\]

The notes interpret the question as the complement
of getting a success in every single trial.

So we calculate:
\[
1-P(Y=16).
\]

Here:
\[
Y\sim \operatorname{Bin}(16,0.8).
\]

Calculate the probability of success all 16 times:
\begin{align*}
P(Y=16)
&=
\binom{16}{16}
(0.8)^{16}(0.2)^0
\\[2mm]
&=
(0.8)^{16}
\\[2mm]
&\approx
0.028147.
\end{align*}

Subtract from \(1\):
\begin{align*}
1-P(Y=16)
&=
1-0.028147
\\[2mm]
&=
0.971853.
\end{align*}

So the probability is approximately:
\[
0.9719.
\]

\subsubsection*{Final answer}

\begin{enumerate}[label=\alph*)]
\item
\[
a=\sqrt{3}-1.
\]

\item
\[
p=P\left(X\le \frac{1}{\sqrt{3}}\right)
=\frac{5}{6}.
\]

\item
Using \(p=0.8\):
\[
P(Y>2)=0.99876864
\approx 0.9988.
\]

\item
With \(n=16\):
\[
1-(0.8)^{16}
=0.971853
\approx 0.9719.
\]
\end{enumerate}

\subsubsection*{What to remember}

For a density:
\[
\int f(x)\,dx=1.
\]

For binomial counting:
\[
Y\sim \operatorname{Bin}(n,p).
\]

\[
P(Y=k)
=
\binom{n}{k}
p^k(1-p)^{n-k}.
\]

For ``more than two'', it is often easier to use:
\[
P(Y>2)=1-P(Y\le 2).
\]

\section*{Uppgift 4}

\subsection*{Original question}

Betrakta normalfördelade stokastiska variabler
\(X_1,X_2,X_3\) med väntevärde \(\mu=1\) och
standardavvikelse \(\sigma=2\). Vi antar dessutom
att \(X_1,X_2,X_3\) är oberoende.

\begin{enumerate}[label=\alph*)]
\item Beräkna \(P(X_1=1)\) och
\[
P(|X_1|\le 1).
\]
\((2p)\)

\item Beräkna
\[
P(\max(X_1,X_2,X_3)\le 1).
\]
\((2p)\)

\item Låt
\[
Y=X_1-2X_2+3X_3.
\]
Bestäm \(E(Y)\) och \(V(Y)\). \((2p)\)
\end{enumerate}

\subsection*{Compiled notes from Yacoub + Obid}
% UNCERTAIN: The supplied note images are not visibly labelled by author,
% so Yacoub's and Obid's notes are compiled together as the provided notes.

\subsubsection*{What we are doing}

We have three independent normal random variables:
\[
X_1,\ X_2,\ X_3.
\]

Each has:
\[
\mu=1,
\qquad
\sigma=2.
\]

The notes also write:
\[
V(X)=\sigma^2=2^2=4.
\]

\subsubsection*{Given information}

\[
X_i\sim N(1,2),
\qquad
i=1,2,3.
\]

\[
E(X_i)=1.
\]

\[
\operatorname{Var}(X_i)=4.
\]

\[
X_1,X_2,X_3
\text{ are independent.}
\]

\subsubsection*{Step-by-step solution}

\textbf{a) Calculate \(P(X_1=1)\)}

The notes say that \(X_1\) follows a continuous
normal distribution.

For a continuous random variable, the probability of
one exact value is zero:
\[
P(X_1=1)=0.
\]

The notes explain this as:
the area under a single point on a curve is zero.

\textbf{a) Calculate \(P(|X_1|\le 1)\)}

The notes rewrite:
\[
|X_1|\le 1
\]
as:
\[
-1\le X_1\le 1.
\]

Now standardize using the \(Z\)-score formula:
\[
Z=\frac{X-\mu}{\sigma}.
\]

For the lower bound \(X=-1\):
\begin{align*}
Z
&=
\frac{-1-1}{2}
\\[2mm]
&=
\frac{-2}{2}
\\[2mm]
&=-1.
\end{align*}

For the upper bound \(X=1\):
\begin{align*}
Z
&=
\frac{1-1}{2}
\\[2mm]
&=
\frac{0}{2}
\\[2mm]
&=0.
\end{align*}

So:
\begin{align*}
P(-1\le X_1\le 1)
&=
P(-1\le Z\le 0).
\end{align*}

Use the \(Z\)-table:
\[
\Phi(0)=0.5.
\]

\[
\Phi(1)=0.8413.
\]

The notes use:
\[
\Phi(-1)=1-\Phi(1).
\]

So:
\begin{align*}
\Phi(-1)
&=
1-0.8413
\\[2mm]
&=
0.1587.
\end{align*}

Therefore:
\begin{align*}
P(-1\le Z\le 0)
&=
\Phi(0)-\Phi(-1)
\\[2mm]
&=
0.5-0.1587
\\[2mm]
&=
0.3413.
\end{align*}

So:
\[
P(|X_1|\le 1)=0.3413.
\]

\textbf{b) Calculate}
\[
P(\max(X_1,X_2,X_3)\le 1).
\]

Betrakta normalfördelade stokastiska variabler
\(X_1,X_2,X_3\) med väntevärde \(\mu=1\) och
standardavvikelse \(\sigma=2\). Vi antar dessutom
att \(X_1,X_2,X_3\) är oberoende.

\begin{enumerate}[label=\alph*)]
\item Beräkna \(P(X_1=1)\) och
\[
P(|X_1|\le 1).
\]
\((2p)\)

\item Beräkna
\[
P(\max(X_1,X_2,X_3)\le 1).
\]
\((2p)\)

\item Låt
\[
Y=X_1-2X_2+3X_3.
\]
Bestäm \(E(Y)\) och \(V(Y)\). \((2p)\)
\end{enumerate}

The notes rewrite the event:
\[
\max(X_1,X_2,X_3)\le 1
\]
as:
\[
X_1\le 1,\quad X_2\le 1,\quad X_3\le 1.
\]

Because the variables are independent, multiply the
individual probabilities:
\begin{align*}
&P(X_1\le 1,\ X_2\le 1,\ X_3\le 1)
\\[2mm]
&=
P(X_1\le 1)
P(X_2\le 1)
P(X_3\le 1).
\end{align*}

Since all variables have the same distribution:
\[
P(X_1\le 1)
=
P(X_2\le 1)
=
P(X_3\le 1).
\]

Standardize:
\begin{align*}
P(X_1\le 1)
&=
P\left(
Z\le
\frac{1-1}{2}
\right)
\\[2mm]
&=
P(Z\le 0)
\\[2mm]
&=
0.5.
\end{align*}

Thus:
\begin{align*}
P(\max(X_1,X_2,X_3)\le 1)
&=
(0.5)^3
\\[2mm]
&=
0.125.
\end{align*}

\textbf{c) Find \(E(Y)\)}

Betrakta normalfördelade stokastiska variabler
\(X_1,X_2,X_3\) med väntevärde \(\mu=1\) och
standardavvikelse \(\sigma=2\). Vi antar dessutom
att \(X_1,X_2,X_3\) är oberoende.

\begin{enumerate}[label=\alph*)]
\item Beräkna \(P(X_1=1)\) och
\[
P(|X_1|\le 1).
\]
\((2p)\)

\item Beräkna
\[
P(\max(X_1,X_2,X_3)\le 1).
\]
\((2p)\)

\item Låt
\[
Y=X_1-2X_2+3X_3.
\]
Bestäm \(E(Y)\) och \(V(Y)\). \((2p)\)
\end{enumerate}

The notes use the formula:
\[
E(aX+bY)=aE(X)+bE(Y).
\]

Here:
\[
Y=X_1-2X_2+3X_3.
\]

So:
\begin{align*}
E(Y)
&=
E(X_1-2X_2+3X_3)
\\[2mm]
&=
E(X_1)-2E(X_2)+3E(X_3).
\end{align*}

Since each mean is \(1\):
\begin{align*}
E(Y)
&=
1-2(1)+3(1)
\\[2mm]
&=
1-2+3
\\[2mm]
&=
2.
\end{align*}

\textbf{c) Find \(V(Y)\)}

Betrakta normalfördelade stokastiska variabler
\(X_1,X_2,X_3\) med väntevärde \(\mu=1\) och
standardavvikelse \(\sigma=2\). Vi antar dessutom
att \(X_1,X_2,X_3\) är oberoende.

\begin{enumerate}[label=\alph*)]
\item Beräkna \(P(X_1=1)\) och
\[
P(|X_1|\le 1).
\]
\((2p)\)

\item Beräkna
\[
P(\max(X_1,X_2,X_3)\le 1).
\]
\((2p)\)

\item Låt
\[
Y=X_1-2X_2+3X_3.
\]
Bestäm \(E(Y)\) och \(V(Y)\). \((2p)\)
\end{enumerate}

The notes say that because the variables are
independent, we can add variances, but constants
attached to variables must be squared:
\[
\operatorname{Var}(aX)=a^2\operatorname{Var}(X).
\]

So:
\begin{align*}
\operatorname{Var}(Y)
&=
\operatorname{Var}(X_1-2X_2+3X_3)
\\[2mm]
&=
1^2\operatorname{Var}(X_1)
+(-2)^2\operatorname{Var}(X_2)
\\
&\quad
+3^2\operatorname{Var}(X_3).
\end{align*}

Since:
\[
\operatorname{Var}(X_i)=4,
\]
we get:
\begin{align*}
\operatorname{Var}(Y)
&=
1^2(4)+(-2)^2(4)+3^2(4)
\\[2mm]
&=
1(4)+4(4)+9(4)
\\[2mm]
&=
4+16+36
\\[2mm]
&=
56.
\end{align*}

\subsubsection*{Final answer}

\begin{enumerate}[label=\alph*)]
\item
\[
P(X_1=1)=0.
\]

\[
P(|X_1|\le 1)=0.3413.
\]

\item
\[
P(\max(X_1,X_2,X_3)\le 1)
=0.125.
\]

\item
\[
E(Y)=2.
\]

\[
V(Y)=56.
\]
\end{enumerate}

\subsubsection*{What to remember}

For a continuous distribution:
\[
P(X=a)=0.
\]

Standardize normal values using:
\[
Z=\frac{X-\mu}{\sigma}.
\]

For independent events:
\[
P(A\cap B\cap C)=P(A)P(B)P(C).
\]

For independent random variables:
\[
\operatorname{Var}(aX+bY)
=
a^2\operatorname{Var}(X)
+
b^2\operatorname{Var}(Y).
\]

\section*{Uppgift 5}

\subsection*{Original question}

Man har två oberoende observationer
\[
x_1=3
\quad\text{och}\quad
x_2=5
\]
av en stokastisk variabel \(X\) med täthetsfunktionen

\[
f_X(x)=
\begin{cases}
\dfrac{x}{a}
e^{-\frac{x^2}{2a}},
& x>0,\\
0,
& x\le 0,
\end{cases}
\]

där \(a>0\). Bestäm ML-skattningen av \(a\).
\((4p)\)

\subsection*{Compiled notes from Yacoub + Obid}
% UNCERTAIN: The supplied note images are not visibly labelled by author,
% so Yacoub's and Obid's notes are compiled together as the provided notes.

\subsubsection*{What we are doing}

We have two independent observations:
\[
x_1=3,
\qquad
x_2=5.
\]

The notes build the likelihood function by
multiplying the two density values.

Then the notes take \(\ln\) of the likelihood to make
the exponential part easier.

\subsubsection*{Given information}

\[
f_X(x)=
\begin{cases}
\dfrac{x}{a}
e^{-\frac{x^2}{2a}},
& x>0,\\
0,
& x\le 0.
\end{cases}
\]

\[
a>0.
\]

\[
x_1=3,
\qquad
x_2=5.
\]

The observations are independent, so:
\[
L(a)=f_X(x_1)\cdot f_X(x_2).
\]

\subsubsection*{Step-by-step solution}

\textbf{1) Build the likelihood}

Since \(x_1>0\) and \(x_2>0\), use the first part of
the density:
\[
f_X(x_i)
=
\frac{x_i}{a}
e^{-\frac{x_i^2}{2a}}.
\]

The notes write:
\begin{align*}
L(a)
&=
f_X(x_1)\cdot f_X(x_2)
\\[2mm]
&=
\left(
\frac{x_1}{a}
e^{-\frac{x_1^2}{2a}}
\right)
\left(
\frac{x_2}{a}
e^{-\frac{x_2^2}{2a}}
\right).
\end{align*}

Multiply the factors:
\begin{align*}
L(a)
&=
\frac{x_1x_2}{a^2}
e^{-\frac{x_1^2}{2a}}
e^{-\frac{x_2^2}{2a}}
\\[2mm]
&=
\frac{x_1x_2}{a^2}
e^{
-\frac{x_1^2+x_2^2}{2a}
}.
\end{align*}

\textbf{English clarification:}
The notes use the rule
\[
a e^b\cdot b e^c
=ab e^{b+c}.
\]
Here that means the exponents are added.

\textbf{2) Take logarithms}

The notes take \(\ln\) on both sides to make the
exponent less annoying:
\[
\ln(L(a))
=
\ln\left(
\frac{x_1x_2}{a^2}
e^{
-\frac{x_1^2+x_2^2}{2a}
}
\right).
\]

Use the log rules from the notes:
\[
\ln\left(\frac{a}{b}\right)
=
\ln(a)-\ln(b),
\]
\[
\ln(a^2)=2\ln(a),
\]
\[
\ln(e^u)=u.
\]

Then:
\begin{align*}
\ln(L(a))
&=
\ln(x_1x_2)
-
\ln(a^2)
-
\frac{x_1^2+x_2^2}{2a}
\\[2mm]
&=
\ln(x_1x_2)
-
2\ln(a)
-
\frac{x_1^2+x_2^2}{2a}.
\end{align*}

\textbf{3) Differentiate with respect to \(a\)}

The notes differentiate the log-likelihood:
\begin{align*}
\frac{d}{da}\ln(L(a))
&=
\frac{d}{da}\ln(x_1x_2)
-
\frac{d}{da}2\ln(a)
\\
&\quad
-
\frac{d}{da}
\left(
\frac{x_1^2+x_2^2}{2a}
\right)
\\[2mm]
&=
0
-
\frac{2}{a}
+
\frac{x_1^2+x_2^2}{2a^2}
\\[2mm]
&=
-\frac{2}{a}
+
\frac{x_1^2+x_2^2}{2a^2}.
\end{align*}

\textbf{4) Set derivative equal to zero}

Set:
\[
-\frac{2}{a}
+
\frac{x_1^2+x_2^2}{2a^2}
=0.
\]

Now solve for \(a\):
\begin{align*}
\frac{x_1^2+x_2^2}{2a^2}
&=
\frac{2}{a}.
\end{align*}

Multiply both sides by \(2a^2\):
\begin{align*}
x_1^2+x_2^2
&=
4a.
\end{align*}

So:
\[
a
=
\frac{x_1^2+x_2^2}{4}.
\]

\textbf{5) Insert the observed values}

The notes insert:
\[
x_1=3,
\qquad
x_2=5.
\]

Therefore:
\begin{align*}
a
&=
\frac{3^2+5^2}{4}
\\[2mm]
&=
\frac{9+25}{4}
\\[2mm]
&=
\frac{34}{4}
\\[2mm]
&=
8.5.
\end{align*}

So the ML estimate is:
\[
\hat{a}_{ML}=8.5.
\]

\subsubsection*{Final answer}

\[
\hat{a}_{ML}
=
\frac{x_1^2+x_2^2}{4}.
\]

With \(x_1=3\) and \(x_2=5\):
\[
\hat{a}_{ML}=8.5.
\]

\subsubsection*{What to remember}

For independent observations:
\[
L(a)=
\prod_i f_X(x_i).
\]

Using \(\ln(L(a))\) makes products become sums and
exponents easier to differentiate.

For ML estimation:
\[
\frac{d}{da}\ln(L(a))=0.
\]

\section*{Uppgift 6}

\subsection*{Original question}

Vid en undersökning av guldhalten (mätt i
ounce/dwt malm, men enheten spelar ingen roll för
lösningen) från två olika brytplatser \(A\)
respektive \(B\) i en gruva i Sydafrika insamlade
man två oberoende stickprov. Efter en logaritmisk
transformering av observationerna erhölls
normalfördelningslika egenskaper. De logaritmerade
värdena var

\[
\begin{array}{c|c}
\text{Brytplats} & \ln(\text{guldhalt})\\
\hline
A(x_i) &
7.03,\ 6.28,\ 5.51,\ 5.89,\ 5.48,\ 5.77,\ 6.62,\ 4.42\\
B(y_i) &
7.19,\ 4.71,\ 5.90,\ 6.21,\ 5.53,\ 5.88,\ 4.47,\ 7.57,\ 6.82,\ 5.62
\end{array}
\]

Vi antar att \(x_i\) är observationer av
\[
X\in N(\mu_1,\sigma)
\]
och \(y_i\) är observationer av
\[
Y\in N(\mu_2,\sigma).
\]

\begin{enumerate}[label=\alph*)]
\item Beräkna ett konfidensintervall med
konfidensgrad 95\% för skillnaden mellan
väntevärdena av de logaritmerade guldhalterna.
\((4.5p)\)

\item Testa hypotesen att väntevärdena är lika med
signifikansnivå 5\%. \((1.5p)\)
\end{enumerate}

\subsection*{Compiled notes from Yacoub + Obid}
% UNCERTAIN: The supplied note images are not visibly labelled by author,
% so Yacoub's and Obid's notes are compiled together as the provided notes.
% UNCERTAIN: Two supplied images about passenger weights and n=100 do not
% match this gold-content question, so they are not used here.

\subsubsection*{What we are doing}

We compare the mean log-gold-content from two
independent mines, \(A\) and \(B\).

The question assumes both groups are normal-like and
have the same unknown standard deviation \(\sigma\).
So the notes use a pooled variance and a two-sample
\(t\)-interval.

\subsubsection*{Given information}

Mine \(A\):
\[
n_A=8.
\]

\[
\sum x_i=47.0.
\]

\[
\bar{x}
=
\frac{47.0}{8}
=5.875.
\]

\[
s_A^2=0.6401.
\]

Mine \(B\):
\[
n_B=10.
\]

\[
\sum y_i=59.9.
\]

\[
\bar{y}
=5.99.
\]

\[
s_B^2=0.9990.
\]

The notes use:
\[
\text{df}=n_A+n_B-2=16.
\]

\subsubsection*{Step-by-step solution}

\textbf{a) 95\% confidence interval for}
\[
\mu_1-\mu_2.
\]

\textbf{Step 1: Calculate pooled variance}

The notes use:
\[
s_p^2
=
\frac{
(n_A-1)s_A^2+(n_B-1)s_B^2
}{
n_A+n_B-2
}.
\]

Substitute the note values:
\begin{align*}
s_p^2
&=
\frac{
7(0.6401)+9(0.9990)
}{
8+10-2
}
\\[2mm]
&=
0.8420.
\end{align*}

So:
\begin{align*}
s_p
&=
\sqrt{0.8420}
\\[2mm]
&=
0.9176.
\end{align*}

\textbf{Step 2: Find the critical \(t\)-value}

For a 95\% confidence interval:
\[
\alpha=0.05.
\]

Two-sided means:
\[
\frac{\alpha}{2}=0.025.
\]

With:
\[
\text{df}=16,
\]
the notes use:
\[
t_{0.025,16}=2.120.
\]

\textbf{Step 3: Calculate margin of error}

The notes use:
\[
ME
=
t\cdot s_p
\sqrt{
\frac{1}{n_A}
+
\frac{1}{n_B}
}.
\]

Substitute:
\begin{align*}
ME
&=
2.120\cdot 0.9176
\sqrt{
\frac{1}{8}
+
\frac{1}{10}
}
\\[2mm]
&=
0.9226.
\end{align*}

\textbf{Step 4: Build the interval}

The difference in sample means is:
\begin{align*}
\bar{x}-\bar{y}
&=
5.875-5.99
\\[2mm]
&=
-0.115.
\end{align*}

The interval is:
\[
(\bar{x}-\bar{y})\pm ME.
\]

So:
\begin{align*}
-0.115\pm 0.9226
&=
[-1.0376,\ 0.8076].
\end{align*}

\textbf{b) Test equal means at 5\% significance}

Vid en undersökning av guldhalten (mätt i
ounce/dwt malm, men enheten spelar ingen roll för
lösningen) från två olika brytplatser \(A\)
respektive \(B\) i en gruva i Sydafrika insamlade
man två oberoende stickprov. Efter en logaritmisk
transformering av observationerna erhölls
normalfördelningslika egenskaper. De logaritmerade
värdena var

\[
\begin{array}{c|c}
\text{Brytplats} & \ln(\text{guldhalt})\\
\hline
A(x_i) &
7.03,\ 6.28,\ 5.51,\ 5.89,\ 5.48,\ 5.77,\ 6.62,\ 4.42\\
B(y_i) &
7.19,\ 4.71,\ 5.90,\ 6.21,\ 5.53,\ 5.88,\ 4.47,\ 7.57,\ 6.82,\ 5.62
\end{array}
\]

Vi antar att \(x_i\) är observationer av
\[
X\in N(\mu_1,\sigma)
\]
och \(y_i\) är observationer av
\[
Y\in N(\mu_2,\sigma).
\]

\begin{enumerate}[label=\alph*)]
\item Beräkna ett konfidensintervall med
konfidensgrad 95\% för skillnaden mellan
väntevärdena av de logaritmerade guldhalterna.
\((4.5p)\)

\item Testa hypotesen att väntevärdena är lika med
signifikansnivå 5\%. \((1.5p)\)
\end{enumerate}

The hypothesis is:
\[
H_0:\mu_1=\mu_2.
\]

Equivalently:
\[
H_0:\mu_1-\mu_2=0.
\]

The alternative is:
\[
H_1:\mu_1\ne \mu_2.
\]

\textbf{English clarification:}
This is the same two-sided comparison as the
confidence interval in part a).

Since the 95\% confidence interval is:
\[
[-1.0376,\ 0.8076],
\]
and \(0\) lies inside this interval, we do not reject
\(H_0\).

So the notes' result means that at the 5\% level,
there is not enough evidence to say that the two
log-mean gold contents are different.

\subsubsection*{Final answer}

\begin{enumerate}[label=\alph*)]
\item
The 95\% confidence interval for
\(\mu_1-\mu_2\) is:
\[
[-1.0376,\ 0.8076].
\]

\item
Since \(0\) is inside the interval, we do not reject:
\[
H_0:\mu_1=\mu_2.
\]

At 5\% significance, the notes do not show a
significant difference between the two means.
\end{enumerate}

\subsubsection*{What to remember}

For two independent normal samples with common
unknown \(\sigma\), use pooled variance:
\[
s_p^2
=
\frac{
(n_A-1)s_A^2+(n_B-1)s_B^2
}{
n_A+n_B-2
}.
\]

The confidence interval for \(\mu_1-\mu_2\) is:
\[
(\bar{x}-\bar{y})
\pm
t\cdot s_p
\sqrt{
\frac{1}{n_A}
+
\frac{1}{n_B}
}.
\]

If \(0\) is inside the confidence interval for
\(\mu_1-\mu_2\), then we do not reject equal means.

\end{document}
